% !TEX root = ../notes_sovereignai.tex
% Chapter 3: Everywhere. Tailscale mesh, the tailnet as the credential.
% Ollama leaves the box here, and nothing else does: the browser chat and
% the family door belong to the chapter after this one, because they are ways
% of reaching it rather than ways of carrying it.
%%%%%%%%%%%%%%%%%%%%%%%%%%%%%%%%%%%%%%%%%%%%%%%%%%%%%%%%%%%%%%%%%%%%%%

\index{Tailscale}\index{API key}\index{tailnet}\index{device identity}
\chapter{Everywhere}
\printchapterversion{03_everywhere}
\newthought{Reachable, Not Exposed.}
\vspace{1.5em}

% ==============================================================================================
% THE WHY
% ==============================================================================================

\section{The Why}\index{remote access}\index{attack surface}

\newthought{Ollama answers on the machine itself} and nowhere else. That is fine
for a terminal on the rig and useless for the rest of life: the laptop on a train,
the phone in another country, the coding agent in your terminal at work. This
chapter carries its API to those devices without carrying it to anyone else.

\marginnote{\newthought{Reach and proof are one job here.} Getting a request from a
foreign network to your card is one problem, and proving the request came from you
is another. A mesh answers both at once, because a device that cannot join cannot
send, which is why nothing in this chapter issues a password.}

\newthought{The danger to respect} is that Ollama has no password and was built to
listen only to its own machine. The lazy fix, binding it to the network directly,
hands your GPU to the internet\index{port forwarding}: an open \texttt{11434} is a model
anyone can run at your electricity bill's expense, and a prompt log they can read.
So nothing is exposed raw. A private mesh carries you to the box instead, and
membership of that mesh is the credential.

\newthought{By the end of this chapter} you will have:
\begin{itemize}
    \item joined the rig to a private mesh and reached it by name from a foreign network;
    \item served Ollama over HTTPS on that mesh, with no port forwarded;
    \item driven it from a terminal coding agent off-site through two environment variables.
\end{itemize}

% ==============================================================================================
% REACH AND AUTH
% ==============================================================================================

\section{Reach and Auth}\index{mesh VPN}\index{tunnel}\index{authentication}

\newthought{Tailscale} builds a private encrypted mesh between your own devices.
Install it on the rig and on your laptop and phone, and they reach each other by a
stable name as if on one LAN, through firewalls and carrier NAT\index{NAT traversal}, with no port
forwarded and nothing published. The tailnet itself is the authentication: only
your enrolled devices\index{tailnet} can connect, so the box is reachable everywhere by you and
invisible to everyone else\cite{tailscaleSelfHostAI}.

\marginnote{\newthought{Mesh versus tunnel.}\index{Cloudflare!Tunnel} A mesh joins devices you
own into one private network, and it is the whole of this chapter. A tunnel
publishes one service to the public web behind a login, which is what a relative
who will not install a VPN needs and what a raw model API must never be.
Publishing waits until there is a page worth publishing.}

\newthought{The mesh carries the key field with it.}\index{API key!placeholder} A client still
sends an authorisation header, and Ollama still ignores it, so the credential
is the device rather than the string. That is a stronger claim than it sounds:
a leaked key is a key someone else can use, and a leaked tailnet name is a name
that resolves for nobody. What it costs is that every client must be a device you
enrolled, which is the point.

\marginnote{\newthought{When the client cannot join.}\index{bearer token}\index{reverse proxy} A build server or a
shared box you do not own cannot run Tailscale, and then the device stops being
the credential. Put a reverse proxy in front of \texttt{11434} that requires a
bearer token\index{coding agent} and forwards only on a match, and the token becomes the
credential for that one client. Every device you do own stays on the mesh.}

% ==============================================================================================
% USE IT ANYWHERE
% ==============================================================================================

\section{Use It Anywhere}\index{tailscale serve}\index{Tailscale}

\marginnote{\newthought{Install once per device.} Tailscale goes on the rig and on
every device you want to reach it from. The phone and laptop apps are a login, not a
config file.}

\newthought{Join the rig to your mesh}, then reach it by name:

\begin{verbatim}
$ curl -fsSL https://tailscale.com/install.sh | sh
$ sudo tailscale up                    # log in once in the browser
$ tailscale serve --bg 11434           # Ollama over HTTPS on the tailnet
\end{verbatim}

\newthought{Reach it from anywhere.} With Tailscale running on your laptop, the rig
answers at its tailnet name over HTTPS, on cellular, in another country, with
nothing exposed to the public internet and no port forwarded on your router.

\newthought{Point a program at it} with two environment variables, which is the
whole of the client setup for anything that speaks the OpenAI shape:

\begin{verbatim}
$ export OPENAI_API_BASE=https://rig.your-tailnet.ts.net/v1
$ export OPENAI_API_KEY=ollama         # the mesh is the credential
$ curl $OPENAI_API_BASE/chat/completions \
    -H "Authorization: Bearer $OPENAI_API_KEY" \
    -d '{"model":"qwen3.8:27b","messages":[{"role":"user","content":"hi"}]}'
\end{verbatim}

\marginnote{\newthought{Coding agents.} Aider, OpenCode, and Cline all take an
OpenAI base URL and key, so the two variables above are the whole setup. Even Claude
Code can be pointed home with \texttt{ANTHROPIC\_BASE\_URL} set to the same host and
\texttt{ANTHROPIC\_AUTH\_TOKEN} set to any string.}

\newthought{What you now own} travels with you. From a laptop on a foreign network,
your terminal coding agent reads the two variables, sends a task to the rig at home,
and edits your code on a model nobody meters.

% ==============================================================================================
% BRIDGE
% ==============================================================================================

\newthought{Ollama is now everywhere you type}, on every device you own and on
no device you do not. What it still lacks is a face: reaching it means a terminal
and a curl, which is no way to hand a model to the rest of the household. The next
chapter gives it a screen anyone can open and a voice that answers out loud,
without a word leaving the building.

\marginnote{\newthought{feedback}}
